\documentclass[10pt,a4paper]{amsart}
\usepackage[margin=19mm]{geometry}
\usepackage{amsmath,amssymb,mathtools}
\usepackage[scr=rsfs,cal=euler]{mathalfa}
\usepackage{xfrac}
\usepackage[unicode,hidelinks,bookmarks=false]{hyperref}
\newcommand{\ie}{i.e.\ }
\newcommand{\cf}{cf.\ }
\newcommand{\resp}{respectively\ }
\newcommand{\Subsection}{\subsection}
\providecommand{\nobreakdash}{\nobreak\hbox{-}}
\setlength{\emergencystretch}{2em}
\multlinegap0pt
\usepackage{xcolor}
\usepackage{float}
\usepackage{caption}
\usepackage[shortlabels]{enumitem}
\usepackage{tikz}
\usepackage{amssymb}
\usepackage{mathtools}
\usepackage{bbm}
\newtheorem{lemma}{Lemma}
\newtheorem{claim}{Claim}
\title{Graphings with few circulations}
\author{G\'abor Kun, L\'aszl\'o M\'arton T\'oth}
\date{Source: arXiv:2512.17071v1 (CC BY 4.0)}
\begin{document}
\maketitle

\noindent\emph{Source and licence.} Graphings with few circulations. Version arXiv:2512.17071v1 (CC BY 4.0), distributed by arXiv under the Creative Commons Attribution 4.0 International licence, http://creativecommons.org/licenses/by/4.0/, as recorded in the arXiv licence metadata for this exact version and shown on the abstract page https://arxiv.org/abs/2512.17071v1. Full licence notice: \texttt{licenses/arxiv-2512.17071-CC-BY-4.0.txt}.\par\smallskip

\noindent\textit{Editorial scope.} Counted unit: the article's lemma lemma:extension (source0.tex lines 337--351). The article's complete original proof of it is delivered with it (source0.tex lines 352--463, one \texttt{proof} environment of 112 lines). No mathematics is written, repaired or re-derived: the statement and the proof are the article's own lines, in the article's own notation, and no retained line is substituted or re-flowed.  Retained uncounted as the source-local context the proof needs: lines 285--287; lines 307--309.  Nothing was omitted from any retained span, so the package carries no omission record.  No retained line contains a citation, so the wrapper carries no bibliography and no bibliography entry.  No retained line contains a figure environment, an image-inclusion command or drawing code, so no image is delivered and the package declares no asset dependency.  Editorial formatting only, in addition: an AMS \texttt{amsart} preamble that restates the article's own declarations and macros used by the retained lines, namely the environments \texttt{lemma}, \texttt{claim} on separate counters, together with the standard packages those lines need.  The retained environments print in this document as Lemma 1, Claim 1 in order of appearance, whereas the article prints the same items with the numbering of its own full document.  The complete native source of the article as distributed in the arXiv e-print for this exact version is bundled unchanged as \texttt{sources/191287/source0.tex}.
\begin{equation} \label{eqn:edge_lift}
(u,v) \in E(G_n) \textrm{, and } \ \mathbf{y}(\beta)- \mathbf{x}(\beta) = \beta(u,v) \quad \forall \beta \in \mathcal{B}_n.
\end{equation}

\begin{lemma}[Deficiency lemma] \label{lemma:deficiency}
    If $\textbf{m}$ has one single tight coordinate, then $\texttt{def}_{\alpha^{n+1}_i}(V_{\mathbf{m}}) =0$ for all $i \in \{1,\ldots, k\}$. 
\end{lemma}

\begin{lemma}[Extension lemma] \label{lemma:extension}
    If $K_n$ is large enough, we can add vertices $U(G_{n+1})$ and edges such that \begin{enumerate}
        \item \label{itm:proportion} the proportion of new vertices $\frac{|U(G_{n+1})|}{|V(G_{n+1})|}$ is less than $\theta_n$;
        
        \item \label{itm:circ} the $\alpha^{n+1}_i$ already defined on $E(V_0)$ can be extended to circulations with values bounded by $M$ and denominator dividing $D$ such that they remain almost orthogonal: 
        
        $\langle \alpha_i^{n+1}, \alpha_j^{n+1} \rangle \leq \frac{1}{2k}-\frac{1}{2k+2n}$ for $i \neq j$; 
        
        \item \label{itm:bounded_potential_gradient} for each $\beta \in \mathcal{B}_n$ the corresponding coordinate function $p^n_{\beta}$ can be extended to a potential $h_{\beta}$ on $V(G_{n+1})$ such that $\|h_{\beta_{n}}\|_{\infty}\leq K_n M_n$ and $\| \nabla h_{\beta} \|_{\infty} \leq M_n$;
        
        \item{$p_0^n$} \label{itm:proper} extends to a homomorphism $f_n: G_{n+1} \to G_n$ such that the preimage of every edge in $G_n$ has the same size.

        \item The graph $G_{n+1}$ is connected and non-bipartite.
    \end{enumerate}
\end{lemma}

\begin{proof}

First, let us reattach all the edges on the boundary of $V_0(G_{n+1})$ that were cut because they leave the $K_n M_n$-by-$K_nM_n$ box, but now we add distinct endvertices for them. That is, for $e=vw \in E(G_n)$ and $\mathbf{x} \in ([K_n M_n]_{D_n})^{\mathcal{B}_n}$ such that $\mathbf{x}(\beta) + \beta(vw) \notin [K_n M_n]_{D_n}$ for at least one coordinate $\beta$, we add a vertex $u_{e,\mathbf{x}}$ and a $G_{n+1}$-edge $\tilde{e}$ between $(v,\mathbf{x})$ and $u_{e,\mathbf{x}}$, and define $\alpha^{n+1}_i(\tilde{e}) = \alpha^n_i(e)$. Now the flow condition is satisfied at all vertices in $V_0(G_{n+1})$, and all new vertices $u_{e,\mathbf{x}}$ have degree 1. We partition the newly introduced vertices according to the $M_n$-by-$M_n$ cube that their unique neighbor belongs to: $U_{0,\mathbf{m}}=\{u_{e,\mathbf{x}} \mid (v,\mathbf{x}) \in V_{\mathbf{m}}\}$.

Let us explain in this paragraph why almost orthogonality will follow from the construction.
We refer to the set of edges introduced here as the \emph{boundary} of $V_0(G_{n+1})$, denoted by $\partial(V_0)$. We partition $\partial(V_0)$ into $\partial^{+}(V_0)$ and $\partial^{-}(V_0)$ as follows. Fix any order on $\mathcal{B}_n$, and for each $\tilde{e}$ let $\beta_{\tilde{e}}$ denote the first coordinate $\beta$ such that $\mathbf{x}(\beta) + \beta(vw) \notin [K_n M_n]_{D_n}$. If $\mathbf{x}(\beta_{\tilde{e}}) + \beta_{\tilde{e}}(vw) \leq 0$, then $\tilde{e}$ belongs to $\partial^{-}(V_0)$, and if $\mathbf{x}(\beta_{\tilde{e}}) + \beta_{\tilde{e}}(vw) > K_n \cdot M_n$, then $\tilde{e}$ belongs to $\partial^{-}(V_0)$. The point is that the original edges $E(V_0)$ together with, say $\partial^{+}(V_0)$ provide a perfect cover of $G_n$, and the $\alpha^{n+1}_i$ are the lift of the $\alpha^{n}_i$ on these edges. So the sum of $\alpha^{n+1}_i\cdot \alpha^{n+1}_j$ over $E(V_0)\cup\partial^{+}(V_0)$ is close to zero. Thus, the almost orthogonality will be ensured as long as the proportion of newly added edges is small. 

The next two claims perform different surgeries on the boundary boxes. We work with boxes of the form $\underline{V}_{\mathbf{m}}=V_{\mathbf{m}} \cap \underline{V}_0({G_{n+1}})$ and their isomorphic copies. The surgery will add path-like gadgets in order to connect vertices with deficit to vertices with sufficit. First, we consider the case when $\mathbf{m}$ has one single tight coordinate. By the Deficiency Lemma \ref{lemma:deficiency} we can match vertices with deficit to those with sufficit within the same box, hence these paths do not need to be long. Hence, the number of additional vertices will be proportional to the total size of these boxes, that is, $O(1/K_n)|V_0(G_{n+1})|$. We will need to be more careful in the proof of Claim \ref{lemma:surgery2}, when we simultaneously treat all the boxes with at least two tight coordinates. We have to allow these paths to have length $O(K_n)$ in order to extend every $p_{\beta}^n$ to a potential with bounded gradient (i.e., not depending on $K_n$). This will be compensated by the fact that the proportion of these boxes is $O(1/K_n^2)$, so the number of added vertices is again $O(1/K_n)|V_0(G_{n+1})|$.

\begin{claim} \label{lemma:surgery}
If $\mathbf{m}$ has one single tight coordinate, then we can add a set of vertices $U_{\mathbf{m}}$ to $V_0({G_n+1})$ and a set of edges with at least one endvertex in $U_{\mathbf{m}}$ such that 

\begin{enumerate}

\item \label{egy}
$|U_{\mathbf{m}}| \leq c|\underline{V}_{\mathbf{m}}|$, where $c$ depends only on $d, k, DM,$ and $G_n$, but not on $K_n$,

\item 
the vertices of $U_{\mathbf{m}}$ are not adjacent to any vertex of 
$V_0(G_{n+1}) \setminus \underline{V}_{\mathbf{m}}$,

\item 
the resulting graph is $d$-regular at every vertex of $\underline{V}_{\mathbf{m}} \cup U_{\mathbf{m}}$,

\item \label{negy}
the $\alpha^{n+1}_i$ extend to $U_{\mathbf{m}}$ such that the flow condition holds at every vertex of $\underline{V}_{\mathbf{m}} \cup U_{\mathbf{m}}$,


\item \label{ot}
$|\alpha^{n+1}_i| \leq M$ on the set of edges spanned by $\underline{V}_{\mathbf{m}} \cup U_{\mathbf{m}}$,

\item
the homomorphism $p_0^n$ extends to a homomorphism $f_n$ on $U_{\mathbf{m}} \cup V_0(G_{n+1})$.


\end{enumerate} 
\end{claim}

\begin{proof}
We have already added $U_{0,\mathbf{m}}$. From each vertex $u\in U_{0,\mathbf{m}}$, build a $(d-1)$-ary tree of depth $\left\lceil \log_{d-1}\big(k(MD)^k +2 \big)\right\rceil$. At the root $u$, each $\alpha^{n+1}_{i}$ defines an inflow, which we will spread out over the tree. We define the $\alpha_i^{n+1}$ as we go down the tree, splitting the values, and eventually subdivide the  vector $(\alpha_1^{n+1}, \ldots, \alpha^{n+1}_k)$ into parts with at most one nonzero coordinate taking value $\pm \frac{1}{D}$ and multiple vertices with zero coordinates only. We do this for every $u$. By Lemma~\ref{lemma:deficiency}, the number of leaves with value $\frac{1}{D}$ and $-\frac{1}{D}$ in the $i$-th coordinate are the same, for every $i$. Therefore, we can choose a matching on the leaves of the trees connecting the $\frac{1}{D}$ values to the $-\frac{1}{D}$ values. We connect every matched pair of leaves by a gadget graph. The rest of the leaves can be matched since $M_n$ and hence the number of leaves is even. We only make sure that there are two leaves in the same tree matched to each other in order to create an odd cycle, hence the resulting graph will not be bipartite. For the natural number $m$, the gadget graph $\texttt{Gad}(m)$ will consist of $m$ consecutive copies of the complete bipartite graph $K_{d-1,d-1}$, joined by $d$ edges bijectively between the appropriate bipartition classes, and connected to two distinguished vertices at the end. Formally, we set $V(\texttt{Gad}(m))=\{u,v\} \cup [2m] \times [d-1]$, and define the set of edges as follows. 
\begin{align*}
    E(\texttt{Gad}(m))=&\bigg\{\big(u,(1,j) \big), \big( v,(2m,j) \big): j \in [d-1] \bigg\} \\
&\cup \bigg\{\big( (2i,j),(2i+1,j) \big): i \in [m-1], j \in [d-1] \bigg\}\\
 &\cup\bigg\{\big( (2i-1,j),(2i,j') \big): i \in [m], j, j' \in [d-1] \bigg\}.
 \end{align*}

For every matched pair $(a,b) $ we add a copy of $\texttt{Gad}(m(a,b))$ to the graph, where $m(a,b)\leq |V(G_n)|$ will be chosen later, and identify $a$ with $u$ and $b$ with $v$. These gadgets will be pairwise disjoint. We send the flow along a path from $a$ to $b$ in the gadget, and define it to be zero on the other edges. Now $\ref{egy}-\ref{ot}$  hold for this construction. 

It remains to extend $p_0^n$ to a homomorphism $f_n$. We can define $f_n$ arbitrarily on the trees rooted at the vertices of $U_{0,\mathbf{m}}$ as long as we map each child to a neighbor of the image of its ancestor. The gadget $\texttt{Gad}(m)$ has a homomorphism into the path of length $2m+1$ mapping $u$ and $v$ to the endvertices. Therefore, we can choose $m(a,b)$ for a matched pair $(a,b)$ such that there is a walk of length $2m(a,b)+1$ connecting $f_{n}(a)$ and $f_n(b)$. Such a path exists, since $G_n$ is connected and not bipartite.
\end{proof}


Now we treat all the other boxes on the boundary together. Let $V_0^+$ denote the union of the boxes $V_{\mathbf{m}}$ with at least two tight coordinates. Set $\underline{V}_0^+=V_0^+ \cap \underline{V}_0(G_{n+1})$.

\begin{claim}\label{lemma:surgery2}
We can add a set of vertices $U_0^+$ to $V_0(G_{n+1})$ and a set of edges with at least one endvertex in $U_0^+$ such that 

\begin{enumerate}

\item 
$|U_0^+| \leq c K_n |\underline{V}_0^+|$, where $c$ depends only on $d, k, DM, G_n$ and $\mathcal{B}_n$, but not on $K_n$,

\item 
the vertices of $U_0^+$ are not adjacent to any vertex of 
$V_0(G_{n+1}) \setminus \underline{V}_0^+$,

\item 
the resulting graph is $d$-regular at every vertex of $\underline{V}_0^+ \cup U_0^+$,

\item
the lift of each $\alpha^n_i$ extends to $U_0^+$ such that the flow condition holds at every vertex of  $\underline{V}_0^+ \cup U_0^+$,

\item
$|\alpha^{n+1}_i| \leq M$ on the set of edges spanned by $\underline{V}_0^+ \cup U_0^+$,

\item
the homomorphism $p_0^n$ extends to a homomorphism on  $\underline{V}_0^+ \cup U_0^+$,


\item \label{extra} for each $\beta \in \mathcal{B}_n$ the corresponding coordinate function $p^n_{\beta}$ can be extended to a potential $h_{\beta}$ on $\underline{V}_0^+ \cup U_0^+$ such that $\|h_{\beta_{n}}\|_{\infty}\leq K_n M_n$ and $\| \nabla h_{\beta} \|_{\infty} \leq M_n$.
\end{enumerate} 
\end{claim}

\begin{proof}
Note that the total deficiency of any $\alpha_i^{n+1}$ on $\underline{V}_0^+$ is again zero. The construction will be similar to that in the previous lemma.  The difference is that in order to satisfy \ref{extra}, we choose $m(a,b)>K_n$. 
This allows us to extend each $p^n_{\beta}$ to every gadget as a linear function of the distance from $a$ such that \ref{extra} holds, since the distance of $a$ and $b$ is at least $K_n$ inside the gadget.
The rest of the proof is the same. 
\end{proof}

We apply Claim \ref{lemma:surgery} to the boxes with exactly one tight coordinate $\mathbf{m}$, consisting of disjoint isomorphic copies of $\underline{V}_{\mathbf{m}}$, and Claim \ref{lemma:surgery2} to the union of the boxes with at least two tight coordinates, consisting of disjoint isomorphic copies of $\underline{V}_0^+$. These allow us to extend each $\alpha_i^{n+1}$ and the homomorphism $p_0^n$ from $V_0(G_{n+1})$, while each $p^n_{\beta}$ is extended to $h_\beta$ on $U_0^+$. We can define $h_\beta$ as a constant on $U_{\mathbf{m}}$ for every $\mathbf{m}$ with one single tight coordinate. 

{\bf Connectivity of $G_{n+1}$.} Now we apply some surgeries to make $G_{n+1}$ connected. After some preparation we will choose pairs of edges $(a,b)$ and $(a',b')$ in different components, whose removal does not disconnect their component, remove them and add the edges $(a,b')$ and $(a',b)$ in order to merge these components. Before making this surgery we will make sure that the circulations $\alpha_i^{n+1}$ take the same values on these edges, that $(a,b$) and $(a',b')$ are mapped to the same edge of $G_n$, and the values of the potentials $h_{\beta}$ almost agree at $a$ and $a'$, and also at $b$ and $b'$. The first two conditions enable us to maintain $\alpha_i^{n+1}$ and $f_n$ during the surgery, whereas the last one will guarantee that $\|\nabla h_{\beta}\|_{\infty} \leq M_n$ still holds. Hence, in preparation, we will replace certain edges $(x,y)$ by a gadget $\texttt{Gad}(m)$, i.e., we remove the edge $(x,y)$, add an isomorphic copy of $\texttt{Gad}(m)$, and connect $x$ to $u$ and $y$ to $v$. We fix an edge $e$ of $G_n$, and choose the size of the gadget $\texttt{Gad}(m)$ such that there is a walk of length $(2m+3)$ connecting $f_n(x)$ and $f_n(y)$ in $G_n$ going through $e$ at least $2|\mathcal{B}_n|$ times, and define the homomorphism $f_n$ on the gadget accordingly. We extend each $h_{\beta}$ to the gadget linearly, hence $\nabla h_{\beta}$ is at most $M_n/(2m+3)$ on the gadget edges. We do this replacement for every edge in $\partial V_0(G_{n+1})$. We can choose $2|\mathcal{B}_n|$ preimages of $e$ in each gadget introduced above, such that these preimages correspond to different steps of the walk on $G_n$. Removing these edges simultaneously does not disconnect the gadget. We associate two such edges in every gadget to each element of $\mathcal{B}_n$: one to increase that coordinate, and one to decrease it.   

Given two boundary edges whose endvertices differ by exactly $\frac{1}{d^2D_n}$ in exactly one coordinate $\beta \in \mathcal{B}_n$, we choose the edges $(a,b), (a',b')$ associated to $\beta$ from their respective gadgets, and perform the cross-surgery described above. (For one of the boundary edges we choose $(a,b)$ to be the gadget-edge that is meant to decrease the $\beta$ coordinate, and for the other edge, we choose $(a',b')$ to be the gadget-edge meant to increase it, depending on which edge had endvertices with higher $\beta$-coordinate.) The value of $h_{\beta}$ differs on $a$ and $a'$ (and also on $b$ and $b'$) by $\frac{1}{d^2D_n}$, so $\nabla h_{\beta}$ can increase by at most this amount on the edges, implying $\nabla h_{\beta} (a,b') \leq \frac{M_n}{2m+3} + \frac{1}{d^2D_n} \leq M_n$, and similarly for $(a',b)$.

Why is this graph connected? First, note that two vertices on the boundary of $V_0(G_{n+1})$ can be connected if they agree in their coordinate corresponding to $V(G_n)$ and have the same tight coordinate, since every coordinate can be changed by $\frac{1}{d^2D_n}$. (Note that the tight coordinate can also be changed, as long as the result of the change is still a vertex on the boundary of $V_0(G_{n+1})$.) Next, we use our technical assumption on $B_n$. Namely, for every $x \in V(G_n)$ there are at least two elements of $\mathcal{B}_n$ that are not zero on edges incident to $x$. Hence, the faces of the boundary corresponding to these coordinates contain vertices whose first coordinate is $x$, moreover, there is such a vertex which is contained in both faces. Consequently, vertices on the boundary of $V_0(G_{n+1})$, whose first coordinates agree are in the same component. 

Now consider two adjacent vertices $u,v \in V(G_n)$. We claim that there are two adjacent vertices $(u, \mathbf{x}),(v,\mathbf{y})$ on the boundary of $V_0(G_{n+1})$. Consequently, all the vertices on the boundary of $V_0(G_{n+1})$ are in the same connected component.   
To find such vertices $(u, \mathbf{x})$ and $(v,\mathbf{y})$, we use again our technical assumption on $\mathcal{B}_n$. For each of $u$ and $v$ there are two distinct elements of $\mathcal{B}_n$, where vertices with first coordinate $x$ and $y$, respectively, appear on both corresponding faces of the boundary. Hence, we can choose two distinct coordinates, $\beta$ and $\beta'$, that have this property for $u$ and $v$, respectively. We will choose $\mathbf{x}(\beta)$ and $\mathbf{y}(\beta')$ such that they respectively ensure that $(u, \mathbf{x})$ and $(v,\mathbf{y})$ are on the boundary of $V_0(G_{n+1})$. We have to be a bit careful: as we want $(u, \mathbf{x})$ and $(v,\mathbf{y})$ to be connected, choosing $\mathbf{x}(\beta)$ prescribes $\mathbf{y}(\beta)$ via (\ref{eqn:edge_lift}). We assumed that the two faces corresponding to $\beta$ both contain vertices with first coordinate $u$. That is, we can choose $\mathbf{x}(\beta)$ sufficiently close to either $0$ or $K_nM_n$ to be on the boundary, and one of these two choices, depending on the sign of $\beta(u,v)$, will result in $0 < \mathbf{y}(\beta) \leq K_nM_n$. We do the same when choosing $\mathbf{y}(\beta')$: we ensure that $(v,\mathbf{y})$ is on the boundary, and $0<\mathbf{x}(\beta')\leq K_nM_n$. We are free to choose the rest of the coordinates of $\mathbf{x}$ and $\mathbf{y}$ arbitrarily as long as their difference satisfies (\ref{eqn:edge_lift}) and all the values are in $[K_nM_n]_{d^2D_n}$.

Finally, we argue that all vertices are connected to the boundary. Indeed, given an arbitrary vertex $(u,\mathbf{x})$ of $V_0(G_{n+1})$, we choose a cycle in $G_n$ containing $u$ such that some $\beta \in \mathcal{B}_n$ does not sum to zero on the cycle. Iteratively following the lifts of this cycle will keep increasing or decreasing the $\beta$ coordinate, and we will eventually hit the boundary. This finishes the proof of the connectivity of $G_{n+1}$.

{\bf Making the sequence proper.} In order to guarantee that the preimage of every edge in $G_n$ has the same size, we iteratively replace edges of $G_{n+1}$ whose image has a smaller preimage than the maximum over $E(G_n)$ by the above gadget $\texttt{Gad}(1)$. We might map all the new edges of $\texttt{Gad}(1)$ to the same edge of $G_n$ as the old one. We increase the size of the preimage of an edge by $d^2$ each step. The size of every preimage is divisible by $d^2$, since we made the same surgery in every copy of $\underline{V}_0(G_{n+1})$, and cross-surgeries also preserve this property. Thus, \ref{itm:proper} will hold in the end of this process. And the size of the largest preimage has not changed, hence the number of vertices not in $V_0(G_{n+1})$ can become at most $|E(G_n)|$ times bigger during this procedure.
Every potential can be extended to these gadgets in order to satisfy \ref{itm:bounded_potential_gradient}.

{\bf Bounding the size of the boundary.} Note that $|V_0^+| \leq {|\mathcal{B}_n| \choose 2} \frac{4}{K_n^2}|V_0(G_{n+1})|$, while the union of the boxes with one single tight coordinate has size at most 
$|\mathcal{B}_n| \frac{2}{K_n}|V_0(G_{n+1})|$.
Consequently, the number of vertices added using Claims~\ref{lemma:surgery} and \ref{lemma:surgery2}, the cross-surgery, and the surgeries making the size of the preimages equal  is  $O(\frac{1}{K_n})|V_0(G_{n+1})|$. Hence \ref{itm:proportion} follows if $K_n$ is large enough.

{\bf Almost orthogonality.} We prove almost orthogonality by induction. The contribution of the original edges $E_0(G_{n+1})$ together with $\partial^+(V_0)$ is bounded by the induction. The circulations take value at most $M$ on the rest of the edges, whose proportion is bounded by $\theta_n$. So we will be able to make sure that the contribution of the latter is at most $\frac{1}{2k}-\frac{1}{2k+2n} - \big(\frac{1}{2k}-\frac{1}{2k+2n-2}\big)=\frac{1}{2(k+n)(k+n-1)}>0$. This holds if $\theta_n \leq \frac{1}{2(k+n)(k+n-1)M^2}$. 
%\frac{\|\alpha_n^i\|^2}{T^2}$.}
\end{proof} 

\end{document}
